\section{第一部分：将非形式化思考融入形式化证明——LeanSTaR}

\subsection{动机：超越纯形式化训练}

传统方法仅在形式化证明数据上训练模型（状态 $\to$ 下一步策略），但这可能不足以学到产生代码所需的\textbf{思维过程}。

\begin{importantbox}{LeanSTaR 的核心思想}
训练模型在每步形式化证明之前，先生成自由形式的\textbf{非形式化思考}（informal thought）：
\begin{itemize}
    \item 模型可以分解问题、规划下一步、执行非形式化计算
    \item 思考帮助多样化搜索空间
    \item 类似 Chain-of-Thought，增强模型的表达能力
\end{itemize}
\end{importantbox}

\subsection{LeanSTaR 算法}

两阶段训练：

\textbf{阶段一：初始化}
\begin{enumerate}
    \item 对已有证明（如 Mathlib），给定状态和下一步策略，让 LLM 回顾性地生成描述该转换的思考
    \item 收集数十万条（状态, 思考, 策略）样本
    \item 训练初始模型学会在策略前生成思考
\end{enumerate}

\textbf{阶段二：强化学习（Expert Iteration）}
\begin{enumerate}
    \item 模型并行生成多个完整证明
    \item Lean 检验证明是否正确
    \item 收集成功的轨迹加入训练集
    \item 在扩展数据集上训练新模型
    \item 迭代此过程
\end{enumerate}

\begin{warningbox}{树搜索在加入思考后的困难}
传统的 Best-First Search 在加入非形式化思考后效果不佳——难以对思考进行准确评分。LeanSTaR 改用简单的\textbf{并行采样}策略：多条路径并行生成，遇错重试，成功即停。
\end{warningbox}

\subsection{实验结果}

\begin{itemize}
    \item LeanSTaR 在 mini-F2F 基准上将模型从低于 GPT-4 Agent 基线提升至超越它
    \item \textbf{关键发现}：带思考的模型从推理预算扩展（scaling inference budget）中获益更大——思考有效多样化了搜索空间
    \item 模型在非形式化思考中执行了形式化代码中不存在的推理步骤
\end{itemize}

此后，非形式化思考 + RL 训练已被定理证明社区广泛采用（DeepSeek-Prover 等），也是 O1、DeepSeek-R1 等推理模型的核心范式。

\subsection{本章小结}

LeanSTaR 的核心洞察：仅在形式化代码上训练可能不足以学到生成代码所需的思维过程。通过 RL 学习生成非形式化思考，可以显著提升证明能力。

